% !TEX root = ./../main.tex

\section{Experiment}
\label{sec:experiment}

This section presents the examples generated by existing long video generation methods including FIFO-Diffusion, and evaluates their performance qualitatively and quantitatively.
We also perform the ablation study to verify the benefit of latent partitioning and lookahead denoising introduced in FIFO-Diffusion.

\subsection{Implementation details}
\label{subsec:implementation}

We implement FIFO-Diffusion based on existing open-source text-to-video diffusion models trained on short video clips, including three U-Net-based models, VideoCrafter1~\citep{chen2023videocrafter1},  VideoCrafter2~\citep{chen2024videocrafter2}, and zeroscope\footnote{\url{https://huggingface.co/cerspense/zeroscope\_v2\_576w}}, as well as a DiT-based model, Open-Sora Plan\footnote{\url{https://github.com/PKU-YuanGroup/Open-Sora-Plan}}.
We employ the DDIM sampling~\citep{song2021denoising} with $\eta\in\{0.5, 1\}$.
\cref{app:implementation_details} provides more details about our implementations.

For quantitative evaluation, we measure $\text{FVD}_{128}$~\citep{unterthiner2018towards} and IS~\citep{salimans2016is} scores using Latte~\citep{ma2024latte} as a base model, which is a DiT-based video model trained on UCF-101~\citep{soomro2012ucf}.
We generate 2,048 videos with 128 frames each to calculate $\text{FVD}_{128}$, and randomly sample a 16-frame clip from each video to measure IS score, following evaluation guidelines in \citep{skorokhodov2022styleganv}. To calculate computational cost, we adopt VideoCrafter2 as the baseline model, using a DDPM scheduler with 64 inference steps on A6000 GPUs.


\subsection{Qualitative results}
\label{subsec:qual}

\begin{figure}[!t]
    \centering
    \renewcommand{\arraystretch}{0.7}
\scalebox{0.98}{
    \setlength{\tabcolsep}{1pt} \hspace{-0.5mm}
        \hspace{-3mm}
        \begin{tabular}{ccccc}
	        \includegraphics[width=0.2\linewidth]{fig/opensora_fifo_DDPM_jpgs/6/0.jpg} &
            \includegraphics[width=0.2\linewidth]{fig/opensora_fifo_DDPM_jpgs/6/4.jpg} &
            \includegraphics[width=0.2\linewidth]{fig/opensora_fifo_DDPM_jpgs/6/8.jpg} &
            \includegraphics[width=0.2\linewidth]{fig/opensora_fifo_DDPM_jpgs/6/12.jpg} &
            \includegraphics[width=0.2\linewidth]{fig/opensora_fifo_DDPM_jpgs/6/16.jpg} \\
            \multicolumn{5}{c}{\small (a) \textsf{"a serene winter scene in a forest. The forest is blanketed in a thick layer of snow, which~$\text{\myldots}$"}}\\
            % 
            \includegraphics[width=0.2\linewidth]{fig/videocrafter2_jpgs/17/0.jpg} &
            \includegraphics[width=0.2\linewidth]{fig/videocrafter2_jpgs/17/3.jpg} &
            \includegraphics[width=0.2\linewidth]{fig/videocrafter2_jpgs/17/6.jpg} &
            \includegraphics[width=0.2\linewidth]{fig/videocrafter2_jpgs/17/9.jpg} &
            \includegraphics[width=0.2\linewidth]{fig/videocrafter2_jpgs/17/12.jpg} \\
            \multicolumn{5}{c}{\small (b) \textsf{"A vibrant underwater scene of a scuba diver exploring a shipwreck, 2K, photorealistic."}} \vspace{1mm}\\
            %
            \iffalse
            \includegraphics[width=0.2\linewidth]{fig/mp_jpgs/6/0.jpg} &
            \includegraphics[width=0.2\linewidth]{fig/mp_jpgs/6/4.jpg} &
            \includegraphics[width=0.2\linewidth]{fig/mp_jpgs/6/8.jpg} &
            \includegraphics[width=0.2\linewidth]{fig/mp_jpgs/6/12.jpg} &
            \includegraphics[width=0.2\linewidth]{fig/mp_jpgs/6/15.jpg} \\
            \multicolumn{5}{c}{\small (c) \textsf{"A tiger \textbf{\textit{walking}} $\rightarrow$ \textbf{\textit{standing}} $\rightarrow$  \textbf{\textit{resting}} on the grassland, photorealistic, 4k, high definition"}} \vspace{1mm}\\
            \fi
%            \iffalse
            \includegraphics[width=0.2\linewidth]{fig/mp_jpgs/6/1.jpg} &
        \includegraphics[width=0.2\linewidth]{fig/mp_jpgs/6/2.jpg} &
        \includegraphics[width=0.2\linewidth]{fig/mp_jpgs/6/3.jpg} &
        \includegraphics[width=0.2\linewidth]{fig/mp_jpgs/6/4.jpg} &
        \includegraphics[width=0.2\linewidth]{fig/mp_jpgs/6/5.jpg} \\
        \includegraphics[width=0.2\linewidth]{fig/mp_jpgs/6/6.jpg} &
        \includegraphics[width=0.2\linewidth]{fig/mp_jpgs/6/7.jpg} &
        \includegraphics[width=0.2\linewidth]{fig/mp_jpgs/6/8.jpg} &
        \includegraphics[width=0.2\linewidth]{fig/mp_jpgs/6/9.jpg} &
         \includegraphics[width=0.2\linewidth]{fig/mp_jpgs/6/10.jpg} \\
        \includegraphics[width=0.2\linewidth]{fig/mp_jpgs/6/11.jpg} &
        \includegraphics[width=0.2\linewidth]{fig/mp_jpgs/6/12.jpg} &
        \includegraphics[width=0.2\linewidth]{fig/mp_jpgs/6/13.jpg} &
        \includegraphics[width=0.2\linewidth]{fig/mp_jpgs/6/14.jpg} &
        \includegraphics[width=0.2\linewidth]{fig/mp_jpgs/6/15.jpg} \\
         \multicolumn{5}{c}{(c) \small \textsf{"A tiger  \textbf{\textit{walking}} $\rightarrow$  \textbf{\textit{standing}} $\rightarrow$  \textbf{\textit{resting}} on the grassland, photorealistic, 4k, high definition"}} \vspace{4pt} \\
%         \fi
        \end{tabular}
    }
        \caption{
            Illustrations of long videos generated by FIFO-Diffusion based on (a) Open-Sora Plan and (b) VideoCrafter2, as well as (c) multiple prompts based on VideoCrafter2.
            The number on the top-left corner of each frame indicates the frame index.
        } \label{fig:qual_short}
%       \vspace{-4mm}
\end{figure}

\begin{figure}[!h]
    \centering
    \renewcommand{\arraystretch}{0.7}
    \scalebox{0.98}{
    \setlength{\tabcolsep}{1pt}
    \hspace{-3mm}
    \begin{tabular}{cccccc}
        % 
    \rotatebox[origin=c]{90}{\small Ours \hspace{-16mm}} &
        \includegraphics[width=0.2\linewidth]{fig/videocrafter2_jpgs/19/0.jpg} &
        \includegraphics[width=0.2\linewidth]{fig/videocrafter2_jpgs/19/3.jpg} &
        \includegraphics[width=0.2\linewidth]{fig/videocrafter2_jpgs/19/6.jpg} &
        \includegraphics[width=0.2\linewidth]{fig/videocrafter2_jpgs/19/9.jpg} &
        \includegraphics[width=0.2\linewidth]{fig/videocrafter2_jpgs/19/12.jpg} \vspace{1mm} \\
    %
        \rotatebox[origin=c]{90}{\small FreeNoise \hspace{-16mm}} &
        \includegraphics[width=0.2\linewidth]{fig/freenoise_jpgs/5/0.jpg} &
        \includegraphics[width=0.2\linewidth]{fig/freenoise_jpgs/5/3.jpg} &
        \includegraphics[width=0.2\linewidth]{fig/freenoise_jpgs/5/6.jpg} &
        \includegraphics[width=0.2\linewidth]{fig/freenoise_jpgs/5/9.jpg} &
        \includegraphics[width=0.2\linewidth]{fig/freenoise_jpgs/5/12.jpg} \vspace{1mm}\\
    %
        \rotatebox[origin=c]{90}{\small Gen-L-Video \hspace{-16mm}} &
        \includegraphics[width=0.2\linewidth]{fig/genlvideo_jpgs/5/0.jpg} &
        \includegraphics[width=0.2\linewidth]{fig/genlvideo_jpgs/5/3.jpg} &
        \includegraphics[width=0.2\linewidth]{fig/genlvideo_jpgs/5/6.jpg} &
        \includegraphics[width=0.2\linewidth]{fig/genlvideo_jpgs/5/9.jpg} &
        \includegraphics[width=0.2\linewidth]{fig/genlvideo_jpgs/5/12.jpg} \vspace{1mm} \\
    %
        \rotatebox[origin=c]{90}{\small LaVie+SEINE \hspace{-16mm}} &
        \includegraphics[width=0.2\linewidth]{fig/lavie_seine_jpgs/18/0.jpg} &
        \includegraphics[width=0.2\linewidth]{fig/lavie_seine_jpgs/18/3.jpg} &
        \includegraphics[width=0.2\linewidth]{fig/lavie_seine_jpgs/18/6.jpg} &
        \includegraphics[width=0.2\linewidth]{fig/lavie_seine_jpgs/18/9.jpg} &
        \includegraphics[width=0.2\linewidth]{fig/lavie_seine_jpgs/18/12.jpg} \vspace{1mm} \\
        & \multicolumn{5}{c}{\small \textsf{"An astronaut floating in space, high quality, 4K resolution."}} \vspace{3pt}\\
    
    \end{tabular}
    }
    \caption{
        Sample videos generated by 
        (first) FIFO-DIffusion on VideoCrafter2, 
        (second) FreeNoise on VideoCrafter2,
        (third) Gen-L-Video on VideoCrafter2, and 
        (last) LaVie + SEINE.
        The number on the top-left corner of each frame indicates the frame index.
    }\label{fig:qual_comparison}
    \vspace{-3mm}
\end{figure} 

We first evaluate the performance of the proposed approach qualitatively.
\cref{fig:qual_long} illustrates examples of long videos~(longer than 10K frames) generated by FIFO-Diffusion based on VideoCrafter2.
It demonstrates the ability of FIFO-Diffusion to generate significantly longer videos than the target length of pretrained baseline models---16 frames in this case.
The individual frames exhibit outstanding visual quality with no perceptual quality degradation even in the later part of the videos while preserving semantic information across all frames.
\cref{fig:qual_short} (a) and (b) present the generated videos with natural motion of scenes and cameras; the consistency of motion is effectively maintained by referencing earlier frames through the generation process.

Furthermore, \cref{fig:qual_short} (c) illustrates that FIFO-Diffusion can generate videos with extensive motion driven by a sequence of changing prompts. 
The capability to generate multiple motions and seamless transitions between scenes highlight the practicality of our method.
Please refer to \cref{app:qual,app:mp} for more examples and our project page\footref{fn:project} for video demos, in comparisons with the videos from other baselines. 
%


\begin{wrapfigure}{r}{0.5\linewidth}
    \centering
%    \vspace{-4mm}
    \includegraphics[width=\linewidth]{./fig/user_study3.pdf} 
    \caption{The results of user study between FIFO-Diffusion and FreeNoise for five criteria.}\label{fig:user_study}
    \vspace{-2mm}
\end{wrapfigure}
%
\cref{fig:qual_comparison} compares the results from FIFO-Diffusion with two training-free techniques, FreeNoise~\citep{qiu2023freenoise} and Gen-L-Video~\citep{wang2023genl} based on VideoCrafter2, as well as a training-based chunked autoregressive method, LaVie~\citep{wang2023lavie} + SEINE~\citep{chen2023seine}.
Note that the chunked autoregressive method requires two models: LaVie for T2V and SEINE for I2V.
We observe that our method significantly outperforms the others in terms of motion smoothness, frame quality, and scene diversity.


Among the training-free methods, Gen-L-Video often produces videos with blurred background while FreeNoise struggles to generate dynamic scenes.\footnote{We provide qauantitative evaluation on the magnitude of motion in \cref{app:motion}.}
The videos from LaVie + SEINE gradually degrade and diverge from text prompts due to error accumulation in their autoregressive generation processes.
Additionally, they often exhibit discontinuities between adjacent chunks, as only the last frame of each chunk is employed to transfer contextual information to the next.
\cref{fig:qual_comparison_app1,fig:qual_comparison_app2} in \cref{app:qual_comparison} provide further examples comparing these methods.

We also conduct a user study to evaluate the long video generation performance of FIFO-Diffusion compared to an existing approach, FreeNoise. 
As shown in \cref{fig:user_study}, users expressed a strong preference for FIFO-Diffusion across all criteria, particularly those related to motion. 
Given that motion is one of the most defining characteristics of videos as opposed to images, the strong performance of FIFO-Diffusion in these criteria is promising and highlights its potential to generate more natural, dynamic videos. 
Details about the user study are provided in \cref{app:user_study_details}.



\begin{table}[t]
    \caption{Comparisons of $\text{FVD}_{128}$ and IS scores on UCF-101. FIFO-Diffusion with latent partitioning and lookahead denoising utilizes Latte~\cite{ma2024latte} as its baseline, where the number of partitions is four ($n=4$).
    The FVD and IS scores of the other algorithms are obtained from their respective papers, and PVDM~\citep{yu2023pvdm} denotes PVDM-L~(400-400s).\vspace{2mm}}
    \centering
    \renewcommand{\arraystretch}{1}
    \setlength{\tabcolsep}{8mm}
    \scalebox{0.8}{
    \begin{tabular}{ccccccc}
        \toprule
        & $\text{FVD}_{128}~(\downarrow)$ & IS~$(\uparrow)$ \\
                \hline
                StyleGAN-V~\citep{skorokhodov2022styleganv} & 1773.4 & 23.94{$\pm$0.73} \\
                VIDM~\citep{mei2022vidm} & 1531.9 & -- \\
                PVDM~\citep{yu2023pvdm} & \ \ 648.4 & 74.40{$\pm$1.25} \\
                %\hline
                %Latte~\citep{ma2024latte} & -- & 73.31 \quad\quad \ \   \\
                %Latte~(64 steps) & -- & 69.42{$\pm$3.90}\\
                \hline
                FIFO-Diffusion (ours) & \ \  \textbf{596.64} & \textbf{74.44}{\textbf{$\pm$1.17}} \\
        \bottomrule
    \end{tabular}
    }
    \label{tab:fvd128}
%    \vspace{-4mm}
\end{table}

\vspace{4mm}
\subsection{Quantitative results}
\label{subsec:auto_eval}


We compare FIFO-Diffusion with the baselines trained for long video generation~\citep{skorokhodov2022styleganv,mei2022vidm,yu2023pvdm} in terms of the $\text{FVD}_{128}$ and IS scores.
As shown in~\cref{tab:fvd128}, our approach outperforms all the compared methods including PVDM-L (400-400s)~\citep{yu2023pvdm}, which employs a chunked autoregressive generation strategy.
Note that PVDM-L (400-400s) iteratively generates 16 frames conditioned on the previous outputs over 400 diffusion steps while our approach only requires 64 inference steps (with lookahead denoising) without need for additional training.

\begin{table}[t]
    \centering
    \caption{Memory usages and inference times of long video generation methods. FIFO-Diffusion utilizes latent partitioning with $n=4$ and lookahead denoising.}
    \vspace{2mm}
    \label{tab:memory_time}
    \renewcommand{\arraystretch}{1.0}
    \setlength{\tabcolsep}{2.5mm}
    \scalebox{0.8}{
        \begin{tabular}{lrccccccc}
            \toprule
            & & \multicolumn{3}{c}{Memory usage~[MB]~()} & & Inference time~[s/frame]~() & \\
            \cmidrule{3-5}
            Method  & & 128 frames & 256 frames & 512 frames & & \\
            \hline
            %
            FreeNoise~\citep{qiu2023freenoise} & & 26,163 & 44,683 & out of memory & & \ \ 6.09 \\
            Gen-L-Video~\citep{wang2023genl} & & 10,913 & 10,937 & 10,965 & & 22.07 \\
            \hline
            %FIFO-Diffusion~(1 GPU) & & 11245 & 11245 & 11245 & 6.20 \\
            %FIFO-Diffusion~(4 GPUs) & & 12210 & 12210 & 12210 & 1.62 \\
            FIFO-Diffusion~(1 GPU) & & 11,245 & 11,245 & 11,245 & & 12.37 \\
            FIFO-Diffusion~(8 GPUs) & & 13,496 & 13,496 & 13,496 & & \ \ 1.84 \\
            \bottomrule
    \end{tabular}}
\end{table}


\subsection{Computational cost}
\label{subsec:memory_time}

To evaluate computational efficiency, we assess memory usage and inference time per frame for training-free, long video generation methods. 
As shown in \cref{tab:memory_time}, FIFO-Diffusion generates videos of arbitrary lengths with a constant memory allocation, while FreeNoise requires memory proportional to the target video length. 
Although Gen-L-Video maintains nearly constant memory usage, it exhibits the slowest inference speed due to redundant computations. 
Notably, FIFO-Diffusion leverages parallel computation; while incorporating lookahead denoising increases computational demand, utilizing multiple GPUs for parallel processing significantly reduces sampling time.
%




\subsection{Ablation study}
\label{subsec:ablation}

We conduct ablation study to analyze the effect of latent partitioning and lookahead denoising on the performance of FIFO-Diffusion.
\cref{fig:ablation} shows that latent partitioning significantly improves both visual quality and temporal consistency of the generated videos.
Moreover, lookahead denoising further refines the quality of generated videos by facilitating temporal coherency and reducing flickering effects.
The videos on our project page\footnote{\url{https://jjihwan.github.io/projects/FIFO-Diffusion}} clearly demonstrate the benefit of FIFO-Diffusion.

Additionally, \cref{tab:ablation} compares the relative MSE loss~(see \cref{eq:mse_loss}) averaged over all time steps across different ablation settings.
The results show that latent partitioning effectively reduces the training-inference gap caused by diagonal denoising as the number of partitions increases.
Furthermore, lookahead denoising enhances the model's noise prediction accuracy, achieving low relative MSE losses (below 1.0) when used in conjunction with latent partitioning.

\begin{figure}[t]
    \centering
    \renewcommand{\arraystretch}{0.7}
    \scalebox{0.98}{
    \setlength{\tabcolsep}{1pt}
    \hspace{-3mm}
    \begin{tabular}{cccccc}
%        \vspace{-3mm}\\
        \rotatebox[origin=c]{90}{\small DD\hspace{-14mm}} &
        \includegraphics[width=0.2\linewidth]{fig/zeroscope_ablation_jpgs/12/dd/0.jpg} &
        \includegraphics[width=0.2\linewidth]{fig/zeroscope_ablation_jpgs/12/dd/1.jpg} &
        \includegraphics[width=0.2\linewidth]{fig/zeroscope_ablation_jpgs/12/dd/2.jpg} &
        \includegraphics[width=0.2\linewidth]{fig/zeroscope_ablation_jpgs/12/dd/3.jpg} &
        \includegraphics[width=0.2\linewidth]{fig/zeroscope_ablation_jpgs/12/dd/4.jpg} \vspace{1mm} \\
        % \shortstack[]{FreeNoise\\\vspace{8mm}} &
        \rotatebox[origin=c]{90}{\small DD+LP\hspace{-14mm}} &
        \includegraphics[width=0.2\linewidth]{fig/zeroscope_ablation_jpgs/12/dd_lp/0.jpg} &
        \includegraphics[width=0.2\linewidth]{fig/zeroscope_ablation_jpgs/12/dd_lp/1.jpg} &
        \includegraphics[width=0.2\linewidth]{fig/zeroscope_ablation_jpgs/12/dd_lp/2.jpg} &
        \includegraphics[width=0.2\linewidth]{fig/zeroscope_ablation_jpgs/12/dd_lp/3.jpg} &
        \includegraphics[width=0.2\linewidth]{fig/zeroscope_ablation_jpgs/12/dd_lp/4.jpg} \vspace{1mm}\\
        % \shortstack[]{FIFO-\\Diffusion\vspace{8mm}} &
        \rotatebox[origin=c]{90}{\small DD+LP+LD\hspace{-14mm}} &
        \includegraphics[width=0.2\linewidth]{fig/zeroscope_ablation_jpgs/12/dd_lp_ld/0.jpg} &
        \includegraphics[width=0.2\linewidth]{fig/zeroscope_ablation_jpgs/12/dd_lp_ld/1.jpg} &
        \includegraphics[width=0.2\linewidth]{fig/zeroscope_ablation_jpgs/12/dd_lp_ld/2.jpg} &
        \includegraphics[width=0.2\linewidth]{fig/zeroscope_ablation_jpgs/12/dd_lp_ld/3.jpg} &
        \includegraphics[width=0.2\linewidth]{fig/zeroscope_ablation_jpgs/12/dd_lp_ld/4.jpg} \vspace{1mm} \\
        % \shortstack[]{FIFO-\\Diffusion\vspace{8mm}} &
        & \multicolumn{5}{c}{\textsf{"A scenic cruise ship journey at sunset, ultra HD."}} \vspace{1mm} \\
    \end{tabular}
    }
    \caption{
        Ablation study.
        DD, LP, and LD signifies diagonal denoising, latent partitioning, and lookahead denoising, respectively.
        The number on the top-left corner of each frame indicates the frame index.
    }\label{fig:ablation}
%    \vspace{-3mm}
\end{figure}

\begin{table}[t]
%\vspace{-5mm}
    \caption{
        Relative MSE losses of ablations.
        `LP' and `LD' denote latent partitioning and lookahead denosing, respectively.
        }
    \label{tab:ablation}
    \vspace{2mm}
    \setlength{\tabcolsep}{5mm}
    \centering
    \scalebox{0.8}{
        \begin{tabular}{cc|ccccc}
            \toprule
             & \# of partitions & without LD & with LD \\
            \hline
            without LP & 1& 1.09 & 1.01 \\
            with LP & 2 & 1.04 & 0.99 \\
            with LP & 4 & 1.02 & \textbf{0.98} \\
            \bottomrule
        \end{tabular}
        }
% \vspace{-4mm}
\end{table}
